\section{Library dan Framework Front-End (Sub-CPMK 2.4)}

\textbf{Library} menyediakan fungsi siap pakai yang dipanggil oleh kode Anda (Anda yang mengontrol alur). \textbf{Framework} mengontrol alur dan memanggil kode Anda (inversion of control). Perbedaan ini sering disederhanakan: ``Anda memanggil library; framework memanggil Anda'' \cite{jsInfo}.

\begin{table}[htbp]
\centering
\begin{tabular}{|P{2.5cm}|P{2cm}|P{3.5cm}|P{4cm}|}
\hline
\textbf{Nama} & \textbf{Jenis} & \textbf{Keunggulan} & \textbf{Kapan Digunakan} \\
\hline
React & Library & Komponen, virtual DOM & SPA, UI kompleks \\
\hline
Vue.js & Framework & Mudah dipelajari, reaktif & SPA, progressive \\
\hline
Angular & Framework & Fitur lengkap; TypeScript & Enterprise, tim besar \\
\hline
jQuery & Library & Manipulasi DOM sederhana & Legacy; hindari untuk baru \\
\hline
Alpine.js & Library & Ringan; di HTML langsung & Interaktivitas ringan \\
\hline
\end{tabular}
\caption{Perbandingan library dan framework front-end populer}
\end{table}

Ekosistem front-end modern didominasi oleh tiga framework/library: \textbf{React} (Facebook, komponen-based, paling populer), \textbf{Vue.js} (progressive, mudah dipelajari), dan \textbf{Angular} (Google, TypeScript, full-featured). Ketiganya menggunakan konsep \textit{component-based architecture}: UI dipecah menjadi komponen kecil yang dapat digunakan ulang, masing-masing dengan template, logika, dan (opsional) style sendiri \cite{mdnJsGuide}.

\begin{konsep}
\textbf{Kuasai JavaScript Dulu.} Sebelum menggunakan framework, pastikan pemahaman JavaScript (DOM, event, async, ES6+) sudah kuat. Framework adalah abstraksi di atas JavaScript vanilla---tanpa fondasi yang kuat, debugging dan pemahaman perilaku framework akan sulit. Framework berubah dan berganti; fundamental JavaScript bertahan \cite{mdnJsGuide}.
\end{konsep}

